\documentclass[11pt,a4paper]{article}
\usepackage[T1]{fontenc}
\usepackage{lmodern,amsmath,amssymb,amsthm,mathtools,microtype}
\usepackage[margin=26mm]{geometry}
\usepackage[colorlinks=true,linkcolor=blue,citecolor=blue,urlcolor=blue]{hyperref}
\newtheorem{theorem}{Theorem}[section]
\newtheorem{proposition}[theorem]{Proposition}
\newtheorem{lemma}[theorem]{Lemma}
\newtheorem{corollary}[theorem]{Corollary}
\theoremstyle{definition}
\newtheorem{definition}[theorem]{Definition}
\newtheorem{remark}[theorem]{Remark}
\newcommand{\Q}{\mathbb Q}
\newcommand{\Z}{\mathbb Z}
\newcommand{\C}{\mathbb C}
\newcommand{\hF}{h_{\mathrm F}}
\newcommand{\hFal}{h_{\mathrm{Fal}}}
\newcommand{\deltaFal}{\delta_{\mathrm{Fal}}}
\newcommand{\gr}{\operatorname{gr}^{\mathrm{can}}}
\newcommand{\Jac}{\operatorname{Jac}}
\newcommand{\vol}{\operatorname{vol}}
\setlength{\emergencystretch}{2em}
\title{Height Bounds for Modular Jacobians and Odd Contact Radicals}
\author{}
\date{9 November 2026}
\begin{document}
\maketitle
\begin{abstract}
We prove an absolute bound $\hF(J_1(N))\ll N^6$ for the stable
Faltings height of the Jacobian of the modular curve $X_1(N)$.
The proof retains the genus dependence in estimates for canonical
Green functions, uses a smooth integral cusp to bound the
self-intersection of the relative dualizing sheaf, and applies an
Arakelov height inequality. A separate comparison bounds the height
of an abelian surface occurring as an isogeny factor by the height of
this Jacobian, with only a logarithmic dependence on the degree of
its field of definition. Applied to the five-contact modular
realization used in the Diophantine construction, this gives
$h(a)\ll M_5(a)^{12}$. All constants are independent of the rational
parameter. The modular realization is stated explicitly as the
arithmetic input; no additional $abc$ or Szpiro hypothesis is used.
\end{abstract}

\section{Heights and the five-contact statement}

For a rational number $a=u/v$ in lowest terms, with $v>0$, set
\[
 h(a)=\log\max(|u|,v).
\]
Thus $h$ is logarithmic height. For an abelian variety $A$ over a
number field, $\hF(A)$ denotes its stable Faltings height with the
original Faltings metric. Stability means that a finite extension
over which $A$ has semistable reduction does not change the value.

Put
\[
 S_* =\{p\text{ prime}:p\le76176\},\qquad
 B_* =\prod_{p\in S_*}p,\qquad
 \mathcal D=\{0,1,2401,-343/32\}.
\]
For $a=u/v\in\Q\setminus\mathcal D$ in lowest terms, define
$M_5(a)$ to be the product of the primes $p\notin S_*$ for which
at least one of
\begin{equation}\label{eq:contact-forms}
 v_p(v),\quad v_p(u),\quad v_p(u-v),\quad
 v_p(u-2401v),\quad v_p(32u+343v)
\end{equation}
is positive and odd. Here $v_p(p)=1$. The first form represents
the contact at infinity. A prime appears in $M_5(a)$ only once.
In particular $M_5(a)$ is a positive squarefree integer.

The following is the precise part of the five-contact construction
used below. It is supplied by the geometric and modular argument in
\cite[proof of Theorem \texttt{o8-height}, equations
\textup{(LR3)--(LR5)}, and proof of Theorem
\texttt{quartic-height19}]{Glazer}.

\begin{definition}[Modular realization input]\label{def:input}
For each $a\in\Q\setminus\mathcal D$, write
$R=B_*M_5(a)$. There is an abelian surface $A$ over a number field
such that
\begin{equation}\label{eq:geometric-comparison}
 h(a)\le2\max(0,\hF(A))+128.
\end{equation}
Moreover, one of the following holds:
\begin{enumerate}
\item $\hF(A)<854R$;
\item there are a number field $K$ of degree $d=[K:\Q]\le2304R$
and an integer $N\le210^{300}R^2$ such that $A_K$ is an
isogeny factor over $K$ of $J_1(N)_K$.
\end{enumerate}
\end{definition}

In the cited construction, the first alternative is the case of
complex multiplication. In the other case the surface is an
isogeny factor of the Jacobian of
$X_1(N)/\ker(\varepsilon)$, where $\varepsilon$ is a nebentypus
character. Pullback of divisors from this quotient to $X_1(N)$
then gives the second alternative in Definition~\ref{def:input}.
No bound on the degree of the newform's coefficient field is part of
this input.

\begin{theorem}\label{thm:contact-height}
Assume the modular realization input of Definition~\ref{def:input}.
There is an absolute constant $C_{12}>0$ such that
\[
 h(a)\le C_{12}M_5(a)^{12}
 \qquad(a\in\Q\setminus\mathcal D).
\]
\end{theorem}

The proof will use a bound for modular Jacobians which is independent
of this arithmetic input.

\begin{theorem}\label{thm:jacobian}
There is an absolute constant $C_J>0$ such that
\[
 \hF(J_1(N))\le C_JN^6
\]
whenever the modular curve $X_1(N)$ has genus at least two.
No squarefreeness condition is imposed on $N$.
\end{theorem}

\section{The genus dependence in the Green function estimate}

Let $X=X_1(N)(\C)$ and $g=g(X)\ge2$.
Then $N\ge5$; the group $\Gamma_1(N)$ contains neither $-I$
nor elliptic elements. Its quotient therefore has the usual
Riemann surface measure, without a central stabilizer factor.
The canonical probability measure of $X$ is
\[
 \mu_{\mathrm{can}}=
 \frac{i}{2g}\sum_{j=1}^g\omega_j\wedge\overline{\omega_j},
\]
where the holomorphic differentials $\omega_j$ are orthonormal
for the inner product $(i/2)\int_X\omega\wedge\overline\omega$.
Write $F_\Gamma$ for the nonnegative density satisfying
\[
 \mu_{\mathrm{can}}=\frac{F_\Gamma}{g}\,\mu_{\mathrm{hyp}}
 \quad\hbox{on }\Gamma_1(N)\backslash\mathfrak H.
\]
The canonical Green function $\gr_X(z,w)$ is normalized to have
mean zero against $\mu_{\mathrm{can}}$ and local singularity
$(2\pi)^{-1}\log|z-w|$.

Fix
\[
 \epsilon=(2+\sqrt3)^{-3/2}.
\]
Let $Y_0$ be the complement of the standard cusp disc with
parameter $\epsilon$ in $\operatorname{SL}_2(\Z)\backslash
\mathfrak H$, and let $Y$ be its inverse image in
$\Gamma_1(N)\backslash\mathfrak H$. The radius of the standard
disc at a cusp of width $m$ is $\exp(-2\pi/(m\epsilon))$ in the
local cusp parameter.

\begin{lemma}\label{lem:green}
There are absolute constants $c_1,c_2>0$ such that
\[
 \sup_YF_\Gamma\le c_1,\qquad
 \sup_{Y\times Y}\gr_X
 \le c_2\left(1+\frac{N}{\sqrt g}\right).
\]
\end{lemma}
\begin{proof}
Section~8.2 of \cite{BruinGreen} gives
\[
 \sup_YF_\Gamma\le25.7,\qquad
 \sup_XF_\Gamma\le\max(25.7,0.0126N^2).
\]
The invariant in equation~(7.1) of that paper is
\[
 \zeta_\Gamma=
 \frac1g\int_XF_\Gamma\,\mu_{\mathrm{can}}
 -\frac1{\vol(\Gamma_1(N)\backslash\mathfrak H)}.
\]
Since $\mu_{\mathrm{can}}$ has mass one,
\begin{equation}\label{eq:zeta-genus}
 \zeta_\Gamma\le\frac{\sup_XF_\Gamma}{g}.
\end{equation}
Theorem~7.1(a) of \cite{BruinGreen}, with the uniform congruence
spectral gap and the constants of its Section~8.2, bounds the
Green function on $Y\times Y$ above by
\[
 15800+2S,\qquad
 S=\sqrt{\left(\frac1{4\eta^2}+4\right)137\zeta_\Gamma},
\]
where $\eta>0$ is an absolute spectral lower bound.
The local kernels added to the Green function in that theorem are
nonnegative: they have the form $L(u)-L(2)$ for $1<u\le2$,
where $L(u)=(4\pi)^{-1}\log((u+1)/(u-1))$ is decreasing.
They can therefore be omitted from the upper bound.

Retaining the factor $1/g$ in \eqref{eq:zeta-genus} gives
$S\ll(1+N)/\sqrt g$. Since $g\ge2$, this proves the lemma.
The constants do not depend on the level or on the number of cusps.
\end{proof}

\section{A smooth integral cusp and the Jacobian height}

For a smooth projective curve $X$ over $\Q$, let $e(X)$ denote
the stable arithmetic self-intersection of the relative dualizing
sheaf, equipped with its Arakelov metric. Explicitly, after passing
to a number field $K$ with a minimal regular semistable model
$\mathcal X/\mathcal O_K$, put
\[
 e(X)=\frac{(\overline\omega_{\mathcal X/\mathcal O_K},
                    \overline\omega_{\mathcal X/\mathcal O_K})}
                   {[K:\Q]}.
\]
Also let $\deltaFal(X)$ be the average of Faltings's delta invariant
over the complex embeddings of $K$. These quantities are independent
of the semistable field of definition.

\begin{lemma}\label{lem:self-intersection}
For $X=X_1(N)$ of genus $g\ge2$,
\[
 e(X)\ll g^2+gN^2+g^{3/2}N,
\]
with an absolute implied constant.
\end{lemma}
\begin{proof}
The Tate curve $\operatorname{Tate}(q^N)$ with its point $q$ of
order $N$ defines an integral cusp $O$ on the coarse modular model
of $X_1(N)$. This cusp is smooth over $\Z$, its width is $N$,
and its cotangent line is freely generated over $\Z$ by $dq$.
These assertions, including their validity when $N$ is not
squarefree, are established in \cite[Chapter~III, Section~5.1,
pp.~96--97]{BruinThesis}.

Let $h_\omega(O)$ be the arithmetic degree of this metrized
cotangent line. With the same normalization of the Green function
as above, the local expression is
\[
 h_\omega(O)=\lim_{z\to O}
 \left(2\pi\gr_X(z,O)-\log|q(z)|\right).
\]
The cusp-disc calculation in the same section of
\cite{BruinThesis} gives
\begin{equation}\label{eq:cusp-height}
 h_\omega(O)\le
 \frac{2\pi}{N\epsilon}
 +\frac{\sup_YF_\Gamma}{g}\frac{N^2\epsilon^2}{4\pi}
 +2\pi\sup_{Y\times Y}\gr_X.
\end{equation}
Lemma~\ref{lem:green} therefore implies
\[
 h_\omega(O)\ll1+\frac{N^2}{g}+\frac{N}{\sqrt g}.
\]

After semistable base change, the arithmetic Hodge index theorem
gives
\[
 e(X)\le4g(g-1)h_\omega(O).
\]
This is the last displayed inequality of
\cite[Chapter~III, Section~5.1, p.~97]{BruinThesis}.
It also follows from \cite[Theorem~2.4.1]{Javanpeykar}:
the stable minimal-model point height is at most the height
computed at the smooth integral cusp. Resolving and then
contracting a regular model does not increase that point height;
see \cite[Lemma~2.3.1]{Javanpeykar}.
Multiplying the preceding height estimate by $4g(g-1)$ proves
the assertion.
\end{proof}

\begin{proof}[Proof of Theorem~\ref{thm:jacobian}]
Theorem~5.6 and Remark~5.8 of \cite{JorgensonKramer} give
\begin{equation}\label{eq:delta}
 \deltaFal(X_1(N))\ll g
\end{equation}
with an absolute implied constant. The result holds for all
levels of genus at least two, after enlarging the constant to
include finitely many small levels. The squarefree-level
hypothesis in Section~6.1 of that paper belongs to a subsequent
arithmetic asymptotic and is not used in \eqref{eq:delta}.

Proposition~2.4.8 of \cite{Javanpeykar} states
\[
 \hFal(X)\le
 \frac{(2g-1)(g+1)}{4(g-1)}e(X)
 +\frac14\deltaFal(X)+20g^3.
\]
Here $\hFal(X)$ is the stable Faltings height of the Jacobian,
as explained in Section~2.3 of that paper. A change of the
constant convention for the metric adds only an absolute
multiple of $g$, which has no effect on the estimate below.
For $g\ge2$ the coefficient of $e(X)$ is at most $2g$.
By Lemma~\ref{lem:self-intersection} and \eqref{eq:delta},
\begin{equation}\label{eq:jacobian-genus}
 \hF(J_1(N))\ll
 g^3+g^2N^2+g^{5/2}N.
\end{equation}
The index formula for $\Gamma_1(N)$ and the usual genus formula
give $g\le N^2$. Each of the three terms on the right of
\eqref{eq:jacobian-genus} is consequently at most $N^6$.
\end{proof}

\section{Comparison with an abelian surface}

The next lemma keeps the dimension of the isogeny comparison
equal to two.

\begin{lemma}\label{lem:surface}
Let $J$ be an abelian variety of dimension $g\ge2$ over a number
field $K$, and let $A$ be an abelian surface which is an isogeny
factor of $J$ over $K$. Put $d=[K:\Q]$ and
\[
 H=\max\left(1,\hF(J)+\frac32g\right).
\]
Then
\[
 \hF(A)\le H+
 8\left(32\log14+\log d+
 \log\max(1,\log d,H)\right).
\]
\end{lemma}
\begin{proof}
Choose a homomorphism $A\to J$ over $K$ with finite kernel and
let $A_0\subset J$ be its image. Thus $A_0$ is an abelian surface
over $K$ and is isogenous over $K$ to $A$.
R\'emond's inequality \cite[Corollary~1.2]{Remond} gives
\[
 \hF(A_0)+\hF(J/A_0)\le\hF(J).
\]
Although that corollary is initially stated for relative heights,
it also holds for stable heights, by making a common semistable
base change; this is explicitly observed after the corollaries
in \cite{Remond}. Bost's lower bound, in the original
normalization, is
\[
 \hF(B)\ge-\log(\pi\sqrt2)\dim B>-\frac32\dim B.
\]
It follows that $\hF(A_0)\le H$.

Apply \cite[Theorem~1.9(1)]{GaudronRemond} with source $A_0$
and target $A$. The theorem provides an isogeny of degree at
most
\[
 \left(14^{32}d\max(1,\log d,\hF(A_0))\right)^{16}.
\]
The increase of stable Faltings height under an isogeny of
degree $D$ is at most $\tfrac12\log D$; see also
\cite[Proposition~2.1]{Remond}. Replacing $\hF(A_0)$ by $H$
in this upper bound proves the result.
\end{proof}

\begin{proof}[Proof of Theorem~\ref{thm:contact-height}]
Fix $a\in\Q\setminus\mathcal D$ and put $R=B_*M_5(a)$.
In the first alternative of Definition~\ref{def:input}, the
comparison \eqref{eq:geometric-comparison} immediately gives
$h(a)\ll R$, which is stronger than required.

In the second alternative, take $K$ and $N$ from that definition.
Since $A$ is a surface occurring as an isogeny factor, the
Jacobian $J_1(N)$ has genus at least two. Its height and
dimension are bounded by $C_JN^6$ and $N^2$, respectively.
Lemma~\ref{lem:surface} therefore yields
\[
 \hF(A)\ll N^6+1+\log d+\log(1+\log d).
\]
The bounds $d\le2304R$ and $N\le210^{300}R^2$ imply
$\hF(A)\ll R^{12}$. Substitution into
\eqref{eq:geometric-comparison} gives $h(a)\ll R^{12}$.
Finally, $B_*$ is fixed and $M_5(a)\ge1$, so the factor
$B_*^{12}$ and the additive constant can be absorbed into a
single absolute constant $C_{12}$.
\end{proof}

\section{Disjoint contact sets and numerical constants}

\begin{corollary}\label{cor:amplification}
Fix a positive integer $r$. For
\[
 \Theta_r=\bigcup_{j=0}^{r-1}
 \left(3j+\left\{0,-\frac12,-1,
                       \frac1{2399},-\frac{32}{407}\right\}\right)
\]
there is a fixed finite exceptional set of primes $\Sigma_r$
and a constant $C_r>0$ such that
\[
 h(t)\le C_r M_r(t)^{12/r}
 \qquad(t\in\Q\setminus\Theta_r).
\]
Here $M_r(t)$ is the product of primes outside $\Sigma_r$
at which some $t-\theta$, $\theta\in\Theta_r$, has positive
odd valuation.
\end{corollary}
\begin{proof}
Apply Theorem~\ref{thm:contact-height} to
$\phi_j(t)=2+1/(t-3j)$ for $0\le j<r$.
The matrix of $\phi_j$ has determinant $-1$, so
$h(t)\le h(\phi_j(t))+O_r(1)$.
Remove all primes dividing coefficients or pairwise determinants
of the $5r$ primitive contact forms, as well as those in $S_*$.
Outside this fixed exceptional set, the radicals from different
blocks have disjoint prime support. Multiplying the $r$ height
inequalities then gives exponent $12/r$. All forms and the
exceptional set are given explicitly in
\cite[Sections~1--2]{Amplification}; its proof applies with
exponent $12$ throughout.
\end{proof}

The constants in Theorems~\ref{thm:contact-height} and
\ref{thm:jacobian} are absolute existence constants. In particular,
the constant in \eqref{eq:delta} has not been evaluated numerically
here. These theorems therefore give no new certified numerical
coefficient height by themselves. They do give uniform exponents
for arguments in which a sufficiently large fixed integer is
already allowed. Any explicit numerical specialization must
add a certified value of the required analytic constants.

\begin{thebibliography}{9}
\bibitem{BruinThesis}
P.~Bruin, \emph{Modular curves, Arakelov theory, algorithmic
applications}, Ph.D. thesis, Universiteit Leiden, 2010,
Chapter~III, Section~5.1, pp.~96--97,
\url{https://math.leidenuniv.nl/scripties/proefschrift-peterbruin.pdf}.

\bibitem{BruinGreen}
P.~Bruin, \emph{Explicit bounds on automorphic and canonical
Green functions of Fuchsian groups}, Mathematika \textbf{60}
(2014), 257--306. Author version dated 4 February 2013,
equation~(7.1), Theorem~7.1(a), and Section~8.2,
\url{https://pub.math.leidenuniv.nl/~bruinpj/green.pdf}.

\bibitem{Javanpeykar}
A.~Javanpeykar, with an appendix by P.~Bruin,
\emph{Polynomial bounds for Arakelov invariants of Belyi curves},
Algebra \& Number Theory \textbf{8} (2014), 89--140,
\url{https://doi.org/10.2140/ant.2014.8.89}.

\bibitem{JorgensonKramer}
J.~Jorgenson and J.~Kramer,
\emph{Bounds on Faltings's delta function through covers},
Annals of Mathematics \textbf{170} (2009), 1--43,
Theorem~5.6 and Remark~5.8,
\url{https://doi.org/10.4007/annals.2009.170.1}.

\bibitem{Remond}
G.~R\'emond, \emph{Propri\'et\'es de la hauteur de Faltings},
Annali della Scuola Normale Superiore di Pisa, Classe di Scienze
\textbf{23} (2022), 1589--1596,
\url{https://doi.org/10.2422/2036-2145.202010_062}.

\bibitem{GaudronRemond}
\'E.~Gaudron and G.~R\'emond,
\emph{Nouveaux th\'eor\`emes d'isog\'enie},
M\'emoires de la Soci\'et\'e Math\'ematique de France
\textbf{176} (2023), Theorem~1.9(1),
\url{https://doi.org/10.24033/msmf.484}.

\bibitem{Glazer}
E.~Glazer and contributors,
\emph{Diophantine definition of the integers in the rationals},
\texttt{manuscripts/quartic\_v9.tex}, commit
\texttt{14e8f601fe0d9173d4a5f084aedc3e07066a7688},
\url{https://github.com/elliotglazer/diophantine-z-q}.

\bibitem{Amplification}
J.~Brox and contributors,
\emph{Height bounds from disjoint sets of rational contacts},
\texttt{contact\_height\_amplification\_20261109\_v13.tex},
\url{https://github.com/JoseBrox/diophantine-z-q_improvements}.
\end{thebibliography}
\end{document}
